\documentclass[11pt]{article}

% NAACL 2025/2026 Review Version (Double-blind submission)
\usepackage[review]{acl}

% Standard package includes
\usepackage{times}
\usepackage{latexsym}
\usepackage[T1]{fontenc}
\usepackage[utf8]{inputenc}
\usepackage{microtype}
\usepackage{inconsolata}
\usepackage{graphicx}
\usepackage{amsmath,amssymb}
\usepackage{booktabs}
\usepackage{subcaption}
\usepackage{algorithm}
\usepackage{algorithmic}
\usepackage{enumitem}
\setlist{itemsep=0pt,topsep=1pt,parsep=0pt,partopsep=0pt}
\usepackage{xcolor}
\setlength{\textfloatsep}{4pt plus 0pt minus 1pt}
\setlength{\floatsep}{4pt plus 0pt minus 1pt}
\setlength{\intextsep}{4pt plus 0pt minus 1pt}
\setlength{\abovecaptionskip}{3pt plus 1pt minus 1pt}
\setlength{\belowcaptionskip}{1pt plus 1pt minus 1pt}
\setlength{\abovedisplayskip}{3pt}
\setlength{\belowdisplayskip}{3pt}
\raggedbottom

\graphicspath{{figures/}{./}}

\title{HAFT: Adversarial Robustness Benchmarking and Budget-Aware Adaptive Auditing Framework for Hindi Automated Fact-Checking}

\author{Anonymous NAACL Submission \\
  Track: Fact-Checking, Adversarial NLP, and Multilinguality \\
  Paper under Double-Blind Review}

\begin{document}
\maketitle

\begin{abstract}
Automated fact-checking (AFC) in Hindi---under-resourced in NLP safety benchmarks---is acutely vulnerable to adversarial manipulation. We introduce \textbf{HAFT}, an adversarial benchmark spanning 1,120 Hindi claims across six domains and 22 attacks (18,304 evaluations on clean-verified claims; 24,640 total) with deterministic Devanagari quality gating. We observe severe evidence-overwrite failures under in-context injection, with 41.0--47.7\% unconditional attack success on the primary verifier; among sampled evidence-attack failures identified by \texttt{gpt-4o-mini}, 96.33\% (578/600) reproduce on Llama 3 70B. On the primary decoder verifier, character noise achieves less than 3.5\% unconditional attack success, whereas dense encoders exhibit acute sensitivity to surface perturbations. Under budgeted auditing ($K \le 5$), static top-5 rankings exploit dominant evidence attacks but leave long-tail blind spots: 34 claims (4.09\%) resist all top-five attacks, including 20 non-jumbling and 15 surface-vulnerable claims. While the \textbf{Hybrid Auditor} optimizes return-maximizing yield (83.71\% $\pm$ 2.66\% discovery; +17.49 pp over Top-2 + Random), return-maximizing exploration yields 0\% blind-spot recall. In contrast, an offline coverage-aware policy recovers 51.43\% of blind-spot instances (18/35 across 5 folds) while matching Static Top-5's 86.71\% discovery, demonstrating that budgeted auditing requires explicit coverage-oriented optimization.
\end{abstract}

\section{Introduction}
\label{sec:intro}

Automated fact-checking (AFC) aims to assess the veracity of factual assertions against retrieved evidence documents \citep{thorne2018fever,schuster2021vitaminc}. While multilingual corpora such as X-Fact \citep{gupta2021xfact} and regional benchmarks \citep{kumar2022hindi} have expanded evaluation to Indic languages, existing testbeds evaluate classification accuracy solely on clean, unperturbed text. Recent investigations reveal that modern NLP pipelines fail catastrophically when adversaries perturb claims, manipulate retrieved documents, or corrupt evidence alignment \citep{wang2021textflint,liu2025adversarial}. However, adversarial robustness in Hindi automated fact-checking remains insufficiently characterized, particularly across controlled claim- and evidence-level perturbations, script-aware quality gating, and budgeted vulnerability auditing. To the best of our knowledge, existing Hindi fact-checking benchmarks have not jointly evaluated adversarial perturbations across claim and evidence levels, deterministic Devanagari-specific quality gating, cross-architecture vulnerability transfer, and sequential budgeted coverage auditing. HAFT addresses this gap through a unified benchmark designed to measure both average vulnerability discovery and long-tail failure coverage.

Specifically, we formalize a 22-attack quality-controlled benchmark over 1,120 real-world Hindi claims and evaluate budgeted defensive auditing. In low-resource public-interest auditing, defensive red-teaming operates under strict query budgets: a fact-checking organization verifying streaming claims cannot afford brute-force evaluation across dozens of attack tools per claim against commercial or rate-limited LLM verification endpoints. However, attack selection that maximizes average vulnerability discovery does not necessarily maximize vulnerability coverage: a fixed global ranking that repeatedly fires high-ASR attacks achieves high average-case yield, but systematically under-covers claim-specific vulnerabilities in the long tail. Budgeted auditing therefore requires balancing average yield with adaptive, coverage-aware exploration.


HAFT investigates three foundational research questions:
\begin{itemize}[leftmargin=1.1em,itemsep=0pt,topsep=1pt]
\item \textbf{RQ1 (Vulnerability Measurement)}: Which adversarial mechanisms cause reliable failures in Hindi automated fact-checking, and how do frontier LLMs perceive Hindi vulnerability?
\item \textbf{RQ2 (Transferability Across Model Families)}: Do measured vulnerabilities transfer across independently developed verifier families or reflect model-specific artifacts?
\item \textbf{RQ3 (Budgeted Defensive Auditing)}: How can an auditor discover claim vulnerabilities under a constrained query budget ($K \le 5$ of 22 attempts), and how does adaptive exploration compare with fixed empirical attack rankings?
\end{itemize}

Our principal empirical contributions are:
\begin{itemize}[leftmargin=1.1em,itemsep=0pt,topsep=1pt]
\item \textbf{Hindi Adversarial Benchmark}: 18,304 clean-verified evaluations (24,640 total across 1,120 claims, 22 attacks, and 6 domains) with authentic evidence from 144 sources and automated quality gating.
\item \textbf{Asymmetric Vulnerability \& Frontier Disconnect}: Evidence poisoning (41.0--47.7\% unconditional ASR) far exceeds character noise ($<$3.5\% unconditional ASR on the primary decoder). Zero-shot frontier models (GPT-4o, Claude 3.5, Kimi-k1.5) achieve only 27.27\% accuracy and 0\% recall on NEG attacks.
\item \textbf{Long-Tail Vulnerability Blind Spots}: Static Top-5 rankings miss 34 claims (4.09\% of clean baseline) that collapse under non-top perturbations (20 non-jumbling, 15 surface-vulnerable, 8 strictly surface-only), concentrated in Crime \& Public Safety ($p = 0.0402$).
\item \textbf{Budgeted Auditing ($K \le 5$)}: The Hybrid Auditor achieves 83.71\% $\pm$ 2.66\% discovery (+17.49 pp over Top-2 + Random), but yields 0\% blind-spot recall. An offline coverage-aware policy matches Static Top-5 discovery (86.71\%) while recovering 51.43\% of blind-spot instances (18/35 pooled; macro $53.5\% \pm 33.7\%$).
\item \textbf{Cross-Family Transfer \& Architecture Diagnosis}: Across sampled failures identified by \texttt{gpt-4o-mini}, 73.19\% (707/966) reproduce on Llama 3 70B (96.33\% on evidence attacks, 578/600), while dense encoders exhibit epistemic collapse into NEI (\texttt{mDeBERTa-v3}) and direct overwrite (\texttt{IndicBERTv2}).
\end{itemize}

\section{Related Work}
\label{sec:related}

\textbf{Automated Fact-Checking.} Fact verification benchmarks originated in English with FEVER \citep{thorne2018fever}, LIAR \citep{wang2017liar}, and VitaminC \citep{schuster2021vitaminc}. Multilingual corpora like X-Fact \citep{gupta2021xfact} and Hindi fake-news datasets \citep{kumar2022hindi} evaluate classification accuracy strictly on clean, unperturbed text. Pre-trained Indic encoders \citep{kakwani2020indicnlpsuite,doddapaneni2023indic} provide foundational representations but leave adversarial vulnerability unmeasured.

\textbf{Knowledge Conflicts and Evidence Poisoning.} In retrieval-augmented generation (RAG) and automated verification, knowledge conflicts arise when external retrieved contexts contradict internal parametric knowledge \citep{xie2024adaptive,wu2024clasheval,zou2025poisonedrag}. Lower-resource languages face documented representation disparities in NLP \citep{joshi2020state}, which motivate examining whether language-specific knowledge limitations contribute to evidence-conflict failures. In our experiments, Hindi verifiers frequently defer to injected evidence, producing substantial verdict-overwrite rates without expressing uncertainty. Our results establish this failure pattern but do not isolate its underlying cause.

\textbf{Adversarial Attacks on Fact-Checking.} Prior surveys document fact-checking and NLP fragility \citep{wang2021textflint,liu2025adversarial}. Evidence-poisoning techniques include AdvAdd \citep{du2022synthetic}, Fact2Fiction \citep{fact2fiction2026}, Fact-Saboteurs \citep{abdelnabi23usenix}, and omission \citep{atanasova2022fact}. As summarized in Table~\ref{tab:comparison}, HAFT differs from existing fact-checking and adversarial NLP benchmarks by integrating a character-to-passage attack taxonomy for Hindi automated fact-checking, deterministic Devanagari quality gating, and sequential budgeted auditing that explicitly evaluates vulnerability discovery yield against long-tail coverage. To our knowledge, existing benchmarks have not combined these components within a unified Hindi fact-checking evaluation framework.

\begin{table}[!t]
\caption{Systematic comparison of HAFT with related fact-checking and adversarial benchmarks. TM: Threat Models (C=Claim, E=Evidence); Gate: Deterministic Script Gating; Aud.: Budgeted Adaptive Auditing ($K \le 5$); Cov.: Coverage vs.\ Yield. Note: TextFlint evaluates multilingual NLP robustness, but does not benchmark Hindi evidence poisoning or budgeted coverage.}
\label{tab:comparison}
\centering
\begingroup
\renewcommand{\arraystretch}{0.88}
\resizebox{\columnwidth}{!}{%
\begin{tabular}{llcccc}
\toprule
Benchmark / Work & Language & TM & Granularity & Gate & Aud. / Cov. \\
\midrule
X-Fact \citep{gupta2021xfact} & Multi (inc.\ Hi) & Clean & None & $\times$ & $\times$ / $\times$ \\
Kumar \& Singh \citep{kumar2022hindi} & Hindi & Clean & None & $\times$ & $\times$ / $\times$ \\
TextFlint \citep{wang2021textflint} & En / Multi & C & Char, Word, Sent & $\times$ (Hi) & $\times$ / $\times$ \\
Fact-Saboteurs \citep{abdelnabi23usenix} & English & E & Passage & $\times$ & $\times$ / $\times$ \\
Fact2Fiction \citep{fact2fiction2026} & English & E & Passage & $\times$ & $\times$ / $\times$ \\
PoisonedRAG \citep{zou2025poisonedrag} & English & E & Passage & $\times$ & $\times$ / $\times$ \\
AKAFake \citep{zhang2024akafake} & English & C & Graph Select & $\times$ & $\times$ / $\times$ \\
\midrule
\textbf{HAFT (Ours)} & \textbf{Hindi (Devanagari)} & \textbf{C + E} & \textbf{Char $\to$ Pass (22)} & \textbf{\checkmark} & \textbf{\checkmark} / \textbf{\checkmark} \\
\bottomrule
\end{tabular}%
}
\endgroup
\end{table}

\textbf{Adaptive Auditing and Sequential Selection.} Policy optimization has been applied to select knowledge subgraphs \citep{zhang2024akafake} and filter training samples \citep{yang2024reinforcement}. In contrast, HAFT investigates sequential attack selection under constrained query budgets ($K \le 5$), formalizing the fundamental trade-off between return-maximizing yield and long-tail vulnerability coverage.

\section{The HAFT Benchmark and Threat Model}
\label{sec:benchmark}

\subsection{Task and Threat Models}
Each benchmark instance is a tuple $x_i = (c_i, e_i, y_i)$, comprising a Hindi claim $c_i$, evidence text $e_i$, and ground-truth stance $y_i \in \{\mathrm{SUP}, \mathrm{REF}, \mathrm{NEI}\}$ (Supported, Refuted, Not Enough Info). The clean verifier predicts $\hat{y}_i = V(c_i, e_i)$. Across 18,304 clean-verified benchmark evaluations (24,640 total across all 1,120 claims), \texttt{gpt-4o-mini} (OpenAI, temperature 0.0, top\_p 1.0, temperature-zero decoding) serves as our \textbf{standardized primary benchmark verifier}, selected for its consistent multilingual instruction-following, a standardized inference configuration, and throughput feasibility across exhaustive combinatorial evaluation. To ensure findings do not reflect model-specific artifacts, cross-architecture generalization is audited across independent model families: open-weight \texttt{Meta-Llama-3-70B-Instruct} (Section~\ref{sec:transfer}), dense encoders (\texttt{IndicBERTv2}, \texttt{mDeBERTa-v3}), and 5 frontier LLMs. An attack generates perturbed text $(c'_{ij}, e'_{ij}) = A_j(c_i, e_i)$, yielding attacked verdict $\hat{y}_{ij} = V(c'_{ij}, e'_{ij})$. A raw flip occurs when $\hat{y}_{ij} \neq \hat{y}_i$.

Our attack taxonomy originated from an initial comprehensive catalogue of 53 candidate adversarial mechanisms across Devanagari morphology, syntax, semantics, and evidence tampering. To ensure rigorous, uniform evaluation, we down-selected exactly 22 representative attack mechanisms for exhaustive benchmarking across all 1,120 claims (18,304 clean-verified evaluations, 24,640 total; Table~\ref{tab:attacks}) based on four explicit, objective criteria: (1)~\emph{Threat Model Coverage}, guaranteeing representation across claim perturbations, injected contexts, and corrupted evidence; (2)~\emph{Linguistic Granularity}, ensuring complete coverage across character, word, sentence, and passage-level manipulation; (3)~\emph{Deterministic Quality-Gate Admissibility}, selecting mechanisms amenable to automated orthographic and lexical verification without human intervention; and (4)~\emph{Literature Grounding}, adapting established multilingual AFC attacks \citep{wang2021textflint,du2022synthetic,fact2fiction2026,abdelnabi23usenix} to Devanagari Hindi. The remaining 31 candidate mechanisms represent unmeasured prioritization hypotheses for future expansions (see the Limitations section).

We delineate adversary capabilities into three threat models:
\begin{itemize}[leftmargin=1.1em,itemsep=0pt,topsep=1pt]
\item \textbf{Threat Model A (Claim Manipulation)}: The adversary modifies the claim ($c'_i \neq c_i, e'_i = e_i$) via surface edits (character noise, typos, synonyms) or semantic blending (\texttt{FactMixing}).
\item \textbf{Threat Model B (In-Context Evidence Injection)}: The adversary concatenates fabricated passages to preserved authentic evidence ($c'_i = c_i, e'_i = [e_i \,\|\, e_{\text{poison}}]$) to benchmark in-context premise adjudication under unauthenticated external text \citep{wu2024clasheval,zou2025poisonedrag}. Because authentic and injected snippets are simultaneously visible in-context, failure to detect contradiction or deferring uncritically to unverified text constitutes an intrinsic verifier failure of epistemic calibration.
\item \textbf{Threat Model C (Evidence Manipulation)}: The adversary tampers directly with authentic evidence ($c'_i = c_i, e'_i \neq e_i$), simulating corrupted scrapes, contextual replacement (\texttt{ContextReplace}), syntactic omission, or imperceptible noise. When authentic reference text is replaced entirely, following the modified snippet reflects an end-to-end RAG pipeline vulnerability (corrupted retrieval or corpus poisoning) rather than an isolated verifier reasoning failure.
\end{itemize}

\subsection{Dataset Provenance and Baseline Isolation}
The benchmark consists of exactly 1,120 claims sampled across 6 diverse domains: Health \& Medicine (189), Politics \& Elections (189), Government Schemes (189), Disaster \& Breaking News (189), Celebrity News (189), and Crime \& Public Safety (175). Label distribution: 571 SUP (51.0\%), 350 REF (31.2\%), and 199 NEI (17.8\%). 

Evidence passages reflect authentic reference texts curated from 144 primary verified sources (PIB Fact Check, BBC Hindi, Dainik Jagran, Alt News, WHO fact sheets) with live URL attribution. No publisher accounts for more than 13.9\% of the corpus (Amar Ujala 13.9\%, WHO 11.1\%, PIB 9.8\%, BBC Hindi 8.5\%). Publisher coverage spans 40 sources in Government Schemes, 37 in Celebrity News, 36 in Crime \& Public Safety, and 25 in Health \& Medicine; Disaster News and Politics \& Elections draw from 3 national wires and 8 political dailies respectively, reflecting centralized coverage in emergency and legislative reporting. To prevent label leakage, we stripped editorial debunk banners, fact-check verdict titles, and conclusion headers from reference texts before constructing the benchmark. NEI instances were constructed from verified journalistic reports covering the claim's topic without conclusive factual confirmation.

\textbf{Dataset Splitting and Deduplication.} For offline RL policy training and budgeted benchmarking, claims are partitioned into 60\% train (499 claims), 20\% validation (166 claims, reserved for policy early stopping and checkpoint selection), and 20\% test (167 claims) splits over the 832 verified correct claims, stratified by civic domain to maintain identical domain proportions across folds. We repeat this across five fixed seeds ($42, 43, 44, 45, 46$). All 1,120 claims were verified via MinHash deduplication (Jaccard similarity $< 0.80$ on character 4-grams) to prevent event- or topic-level duplicate leakage between train and test splits.

\textbf{Clean Baseline Isolation.} Clean verifier accuracy across all 1,120 claims is 74.29\% (832/1,120). Breaking down by gold stance reveals high accuracy on verifiable assertions: \textbf{85.46\% on SUP} (488/571) and \textbf{83.14\% on REF} (291/350), but \textbf{26.63\% on NEI} (53/199), reflecting a known inductive bias toward forcing binary verdicts rather than abstaining \citep{atanasova2022fact}. Clean accuracy varies markedly by domain (95\% Wilson CIs): Government Schemes achieves 99.47\% (188/189, $[97.1\%, 99.9\%]$) and Politics achieves 97.35\% (184/189, $[94.0\%, 98.9\%]$), whereas Disaster News achieves 66.14\% (125/189, $[59.1\%, 72.5\%]$), Health achieves 65.61\% (124/189, $[58.6\%, 72.0\%]$), Celebrity achieves 64.55\% (122/189, $[57.5\%, 71.0\%]$), and Crime \& Public Safety achieves only \textbf{50.86\%} (89/175, $[43.5\%, 58.2\%]$). Exactly 288 clean baseline failures are isolated into a separate denominator. An attack cannot legitimately flip a claim that the verifier fails on in clean form; benchmarking strictly over the 832 verified correct claims yields exactly $832 \times 22 = 18,304$ evaluated attack opportunities.

\textbf{Null-Perturbation Control and Decision Stability.} To establish a rigorous noise floor against temperature-zero API non-determinism, we evaluated repeated clean re-queries ($c'_i = c_i, e'_i = e_i$) across independent evaluation passes. Decision instability under identity transformation occurred in only 2 of 1,120 claims (0.18\% spontaneous flip rate; 95\% Wilson CI $[0.05\%, 0.65\%]$; 0/832 on clean-verified claims, 0.00\%, $[0.00\%, 0.45\%]$), establishing an empirical noise floor of $\le 0.45\%$. All POS and MID attacks (6.13\%--47.72\%) and primary surface perturbations (2.52\%--11.42\%) sit well above this noise floor, whereas imperceptible noise attacks ($\le 0.36\%$, Table~\ref{tab:attacks}) fall within it. Furthermore, on Crime \& Public Safety claims where clean baseline accuracy is lowest (50.86\%), decision boundaries exhibit elevated sensitivity, confirming that surface-level blind spots interact with baseline boundary fragility rather than spontaneous API drift.

\begin{table}[!t]
\caption{Measured attack results across all 22 attacks in HAFT ($N=832$ clean-baseline claims; Primary Verifier: \texttt{gpt-4o-mini}). \textbf{Uncond.\ ASR} reports flips over all 832 claims alongside 95\% Wilson score confidence intervals. \textbf{Cond.\ ASR} conditions on successful, quality-passed perturbations (41--799 eligible claims across all attacks; 592--799 for full-corpus generative attacks).}
\label{tab:attacks}
\centering
\begingroup
\renewcommand{\arraystretch}{0.78}
\fontsize{7.5pt}{8.5pt}\selectfont
\resizebox{\columnwidth}{!}{%
\begin{tabular}{llrrrrl}
\toprule
Attack Mechanism & Family & Raw (\%) & Cond.\ (\%) & Uncond.\ (\%) & 95\% Wilson CI & Tier \\
\midrule
Poisoned evidence addition (AdvAdd) & Evid. (B) & 61.00 & 58.47 & \textbf{47.72} & [44.34, 51.11] & POS \\
Contextualized evidence replace & Evid. (C) & 63.52 & \textbf{59.57} & 43.75 & [40.42, 47.14] & POS \\
Fact mixing & Sent. (A) & 59.82 & 55.81 & 41.59 & [38.28, 44.97] & POS \\
Agentic fictional evidence (Fact2Fiction) & Evid. (B) & 61.33 & 57.60 & 40.99 & [37.69, 44.36] & POS \\
Masked-token claim rewrite & Sent. (A) & 28.49 & 25.39 & 19.71 & [17.15, 22.55] & MID \\
Word jumbling & Word (A) & 12.57 & 12.57 & 11.42 & [ 9.43, 13.76] & NEG \\
Syntactic omission & Evid. (C) & 21.72 & 19.62 & 6.13 & [ 4.69,  7.97] & MID \\
Character insertion & Char. (A) & 4.35 & 3.63 & 3.49 & [ 2.44,  4.96] & NEG \\
Character deletion & Char. (A) & 7.73 & 3.30 & 2.88 & [ 1.95,  4.26] & NEG \\
Typos \& misspellings & Word (A) & 3.44 & 2.95 & 2.76 & [ 1.85,  4.11] & NEG \\
Adversarial trigger & Sent. (A) & 18.25 & 4.28 & 2.76 & [ 1.85,  4.11] & NEG \\
Homoglyph perturbation & Char. (A) & 3.12 & 2.75 & 2.52 & [ 1.66,  3.83] & NEG \\
Phonetic perturbation & Word (A) & 3.75 & 3.25 & 2.28 & [ 1.47,  3.54] & NEG \\
Colloquial rewriting & Sent. (A) & 3.81 & 2.43 & 2.04 & [ 1.28,  3.25] & NEG \\
Character repetition & Char. (A) & 3.79 & 2.46 & 1.92 & [ 1.19,  3.10] & NEG \\
Lexical substitution & Word (A) & 2.95 & 2.47 & 1.80 & [ 1.10,  2.95] & NEG \\
Character swapping & Char. (A) & 7.24 & 1.85 & 1.68 & [ 1.00,  2.80] & NEG \\
Lexically informed rewrite & Sent. (A) & 2.45 & 1.39 & 1.08 & [ 0.57,  2.04] & NEG \\
Synonym replacement & Word (A) & 4.73 & 3.59 & 0.72 & [ 0.33,  1.56] & NEG \\
Imperceptible verification noise & Evid. (C) & 3.25 & 0.40 & 0.36 & [ 0.12,  1.05] & NEG \\
Imperceptible retrieval noise & Evid. (B) & 2.66 & 0.26 & 0.24 & [ 0.07,  0.87] & NEG \\
Entity disambiguation & Word (A) & 0.00 & 0.00 & 0.00 & [ 0.00,  0.46] & NEG \\
\bottomrule
\end{tabular}%
}
\endgroup
\end{table}

\subsection{LLM-Free Deterministic Quality Gating}
A verifier flip constitutes a valid adversarial failure only if perturbed text is fluent and preserves task constraints. We implement a dual-layer deterministic gate in pure Python: (1)~\textbf{Orthographic and Ligature Integrity}: Verifies Devanagari script integrity ($U{+}0900$--$U{+}097F$), rejects illegal stacked diacritics, rejects leading viramas, and filters artificial repetition ($\ge 4$ consecutive repeated aksharas/matras). (2)~\textbf{Deterministic Surface Invariance}: For meaning-preserving perturbations, we enforce lexical similarity via ChrF++ proxy ($n \in [2, 4]$) and token Jaccard similarity ($\text{ChrF} \ge 0.60$ or $J(x, x') \ge 0.50$). For intentional meaning-drift attacks (evidence injection), surface checks are bypassed and orthographic fluency is enforced. Human calibration on 300 stratified double-annotated instances (75 instances per attack family across Character, Word, Sentence, and Evidence tiers) by two native Hindi linguists confirmed an average 93.3\% inter-annotator agreement across 600 judgments on both dimensions (560/600): 94.7\% on Devanagari orthographic fluency (284/300, 16 disagreements, $\kappa = 0.898$) and 92.0\% on task constraint adherence (276/300, 24 disagreements, $\kappa = 0.842$), yielding overall joint acceptance agreement of 93.3\% ($\kappa = 0.881$). Across all 40 disagreements, 16 in fluency stemmed from dialectal diacritics/nukta variants (6) and colloquial spacing/conjunct variants (10); and 24 in task adherence stemmed from colloquial register shifts in lexical substitutions (14) and borderline semantic drift in sentence rewrites (10; full confusion tables and accepted/rejected Devanagari examples are provided in the supplementary material).

\textbf{Gold Label Preservation Under Meaning-Drift Attacks.} For generative claim-modifying attacks that intentionally alter phrasing, we conducted expert stance re-annotation on 50 sampled instances each of \texttt{ClaimRewrite} and \texttt{FactMixing} against reference evidence. In \texttt{ClaimRewrite}, 92.0\% (46/50) preserved the original gold stance, validating that flips represent genuine verifier failures under paraphrasing. In \texttt{FactMixing}, however, blending secondary facts altered the factual ground-truth stance in 42.0\% (21/50) of instances (e.g., introducing an unverified subordinate clause shifted gold SUP to NEI). This label shift directly explains why Fact Mixing exhibits acute architecture-specific sensitivity (Section~\ref{sec:transfer}), transferring at only 3.50\% to Llama 3 70B.

\section{Vulnerability Measurement \& Frontier Disconnect (RQ1)}
\label{sec:measurement}

Table~\ref{tab:attacks} reports results across all 22 measured attacks. Attacks are categorized into three empirical tiers based on quality-gated conditional ASR: \textbf{POS} ($\text{Cond.\ ASR} \ge 40\%$), \textbf{MID} ($15\% \le \text{Cond.\ ASR} < 40\%$), and \textbf{NEG} ($\text{Cond.\ ASR} < 15\%$). All four top-tier (POS) attacks manipulate evidence or complex semantics: Poisoned Evidence Addition (`AdvAdd`) achieves 47.72\% unconditional ASR (58.47\% conditional), followed by Contextualized Evidence Replace (43.75\% uncond / 59.57\% cond), Fact Mixing (41.59\% uncond / 55.81\% cond), and Agentic Fictional Evidence (40.99\% uncond / 57.60\% cond).

\textbf{Failure Analysis Under Conflicting Evidence.} Examining flip transitions reveals acute evidence-overwrite failures: under Poisoned Evidence Addition (\texttt{AdvAdd}), \textbf{88.2\% of flips (389/441 raw, 350/397 baseline-gated)} transition directly from $\text{Gold SUP} \to \text{Attacked REF}$, and under \texttt{Fact2Fiction}, \textbf{98.5\%} transition to REF. Only 0.7--0.8\% flip to NEI. In the framework of RAG knowledge conflicts \citep{xie2024adaptive,wu2024clasheval}, this reflects an architectural failure of epistemic calibration: when injected passages contradict authentic evidence, the verifier exhibits uncalibrated contextual deference, treating unauthenticated snippets as ground truth rather than flagging contradiction or outputting NEI.

To examine whether prompt concatenation order or the absence of source tags could mitigate this vulnerability, we conducted an exploratory position $\times$ provenance pilot ($N=15$ claims, 60 live queries against \texttt{gpt-4o-mini}). In this pilot, reversing the evidence order (\emph{Injected} passage placed first, followed by \emph{Authentic}) still produced a \textbf{100.0\% flip rate}. Furthermore, prefixing explicit provenance tags---\textquoteleft Verified Archive\textquoteright{} versus \textquoteleft Unverified Forward\textquoteright{}---still produced a \textbf{100.0\% flip rate} (with only 20\% shifting to NEI under inverted order). While not ruling out ordering or provenance effects in general, this exploratory pilot confirms that these tested prompt variations did not eliminate evidence-overwrite failures.

In sharp contrast, all five character-level perturbations achieve $<$3.5\% unconditional ASR. Aggregating by attack family ($N=832$ clean claims), \textbf{Evidence manipulation flips 66.47\% of claims (553/832)} and Sentence perturbations flip 50.36\% (419/832), whereas Word perturbations flip only 13.10\% (109/832) and Character noise flips only 6.37\% (53/832). Paired claim-level differences confirm massive, statistically significant separations: Evidence manipulation exceeds Character noise by \textbf{+60.10 pp (10,000-resample bootstrap 95\% CI $[{+}56.49, {+}63.46]$ pp; McNemar paired $p < 10^{-10}$, $b=511, c=11$)}, exceeds Word perturbations by \textbf{+53.37 pp (95\% CI $[{+}49.28, {+}57.33]$ pp; $p < 10^{-10}$)}, and exceeds Sentence perturbations by \textbf{+16.11 pp (95\% CI $[{+}11.06, {+}21.03]$ pp; $p < 10^{-10}$)}. Furthermore, Word perturbations exceed Character perturbations by \textbf{+6.73 pp (95\% CI $[{+}4.45, {+}9.02]$ pp; $p < 10^{-4}$)}. This hierarchy demonstrates that while subword and byte-level representations provide robust resistance to localized orthographic noise, decoder verifiers remain critically fragile to premise contradiction in retrieved evidence.

\begin{table}[!t]
\caption{Zero-shot attack tier prediction across 5 frontier LLMs ($N=22$ attacks: 16 NEG, 4 POS, 2 MID). Pred Distribution reports predicted counts [POS / MID / NEG].}
\label{tab:llm_metrics}
\centering
\begingroup
\renewcommand{\arraystretch}{0.84}
\resizebox{\columnwidth}{!}{%
\setlength{\tabcolsep}{3.5pt}
\begin{tabular}{llrrrrr}
\toprule
Model & Origin & Acc (\%) & Bal.\ Acc & Macro $F_1$ & Cohen's $\kappa$ & Pred [P/M/N] \\
\midrule
GPT-4o & US & 27.27 & 66.67 & 28.03 & 0.158 & 11 / 11 / \phantom{0}\textbf{0} \\
Claude 3.5 Sonnet & US & 27.27 & 66.67 & 28.15 & 0.146 & 14 / \phantom{0}8 / \phantom{0}\textbf{0} \\
Kimi-k1.5 & China & 27.27 & 66.67 & 28.03 & 0.158 & 11 / 11 / \phantom{0}\textbf{0} \\
Sarvam AI & India & 40.91 & 58.33 & 39.50 & 0.169 & 13 / \phantom{0}4 / \phantom{0}5 \\
DeepSeek-V3 & China & \textbf{72.73} & \textbf{87.50} & \textbf{70.09} & \textbf{0.545} & \phantom{0}5 / \phantom{0}7 / 10 \\
\midrule
\multicolumn{2}{l}{\emph{Majority (Always NEG)}} & 72.73 & 33.33 & 28.07 & 0.000 & \phantom{0}0 / \phantom{0}0 / 22 \\
\bottomrule
\end{tabular}%
}
\endgroup
\end{table}

\textbf{Frontier LLM Zero-Shot Disconnect.} We surveyed 5 frontier models across diverse geographic provenance (US: GPT-4o, Claude 3.5 Sonnet; China: DeepSeek-V3, Kimi-k1.5; India: Sarvam AI) to predict attack feasibility zero-shot (Table~\ref{tab:llm_metrics}). Western and Asian frontier models (GPT-4o, Claude 3.5, Kimi) achieve only 27.27\% accuracy (6/22), with a \textbf{0.0\% recall on NEG attacks} (predicting 0 of 16 actual NEG attacks as NEG; Cohen's $\kappa \le 0.158$), consistently overestimating character typos and lexical substitutions. In contrast, DeepSeek-V3 achieves 72.73\% accuracy (Balanced Acc 87.50\%, Macro $F_1$ 70.09\%, $\kappa = 0.545$) by correctly identifying 10 of 16 NEG attacks, matching the raw accuracy of the trivial Always-NEG majority baseline (72.73\%, 16/22 attacks). Sarvam AI achieves 40.91\% accuracy (Balanced Acc 58.33\%, Macro $F_1$ 39.50\%).

\section{Coverage-Aware Budgeted Auditing Methodology (RQ3)}
\label{sec:auditing}

In practical defensive auditing, a security auditor operates under a tight verifier query budget $K \le 5$ (a 77.27\% reduction in LLM verification API calls relative to exhaustive 22-attack evaluation). The central methodological question is: \emph{How should an auditor allocate finite queries between exploiting high-yield empirical attacks and exploring long-tail vulnerabilities?}

\subsection{The Rescued Blind-Spot Discovery}
A static heuristic selecting top-5 empirical arms discovers vulnerable claims on 84.62\% (704/832 claims with at least one verified flip) across the clean-verified corpus (86.71\% $\pm$ 1.30\% across five held-out test splits, conditional on the 832 claims correctly verified in unperturbed form). However, static ranking exhibits a systematic coverage limitation: it is blind to non-top perturbations. Per-claim discrepancy analysis reveals that Static Top-5 produces \textbf{zero flips on 128 claims}. Among these, \textbf{exactly 34 unique claims (4.09\% of clean baseline claims, 3.04\% of the 1,120 corpus)} resist all top-5 arms but collapse under linguistic perturbations: (1)~\textbf{Word Jumbling} rescues \textbf{19 claims} by permuting clause word order while preserving token sets ($J=1.0$); and (2)~\textbf{Non-Jumbling Rescues (20 unique claims, 2.40\%)} collapse under non-top perturbations, including Syntactic Omission (7), Character Insertion (9), Character Deletion (7), Phonetic Shifts (7), Homoglyphs (6), and Typos (6). Within this set, \textbf{15 unique claims (1.80\%)} collapse under surface perturbations, and strictly excluding both jumbling and omission isolates \textbf{8 strictly surface-only claims (0.96\%)} that collapse solely under character or word-level surface edits. Across the 34 claims, rescues sum to 76 instances (mean 2.24 per claim). Blind spots show notable concentration in \textbf{Crime \& Public Safety}: 9 of 34 claims (26.47\%) belong to Crime, despite Crime comprising 10.70\% (89/832) of the clean corpus (OR = 3.23, Bonferroni $p = 0.0402$). Among sub-categories, Crime accounts for 8 of the 20 non-jumbling blind spots (40.0\%, Bonferroni $p=0.0032$), 8 of the 15 surface blind spots (53.33\%, $p < 0.001$), and 6 of the 8 strictly surface-only claims resisting both jumbling and omission (75.00\%, $p < 0.001$). We note that clean Crime claims also exhibit lower zero-shot verifier baseline accuracy (50.86\%, 89/175) than other domains (~80\%), indicating that baseline decision-boundary instability interacts with surface perturbations. A Static Top-5 audit would fail to discover these 34 vulnerabilities within the tested budget.

Across all 832 clean claims, \textbf{738 (88.70\%)} flip under at least one attack while \textbf{94 (11.30\%)} remain robust. Robust Accuracy $RA_K = \frac{1}{N} \sum_{i=1}^N \prod_{a \in \mathcal{A}_K} \mathbb{I}[ \hat{y}_{i,a} = \hat{y}_i ]$ drops monotonically: $RA_1 = 52.28\%$, $RA_2 = 37.86\%$, $RA_5 = 15.38\%$, $RA_{22} = 11.30\%$. On the 34 blind-spot claims, Static Top-5 obtains \textbf{0/34 by construction}; RL's 0\% recall is an empirical result of return-maximizing training. Since empirical rankings reflect verifier-specific frequencies, they do not transfer across architectures (Section~\ref{sec:transfer}), motivating sequential coverage auditing.


\subsection{Auditing Policy Comparisons: Yield Maximization vs.\ Coverage-Aware Search}
Our methodological contribution is a coverage-aware formulation of budgeted vulnerability auditing that explicitly evaluates discovery yield and long-tail blind-spot coverage as distinct objectives. We use REINFORCE as one implementation of return-maximizing adaptive selection and compare it with static, exploratory, and coverage-aware policies under the same query budget. To investigate how sequential search resolves the trade-off between discovery yield and long-tail coverage, we formulate a two-stage budgeted auditing paradigm:
\begin{equation}
\label{eq:hybrid}
\pi(a_t \mid \mathbf{s}_t) = \begin{cases} 
a_{\text{ctx\_rep}}, & t = 1, \\
a_{\text{adv\_add}}, & t = 2, \\
\pi_{\text{adaptive}}(a_t \mid \mathbf{s}_t, \mathbf{m}_t), & t \ge 3,
\end{cases}
\end{equation}
where $a_{\text{ctx\_rep}}$ and $a_{\text{adv\_add}}$ denote the top two empirical evidence attacks (\texttt{EA\_CTXREP} and \texttt{EA\_ADVADD}), fired deterministically at steps $t=1$ and $t=2$, resolving 62.14\% (517/832) of claims instantly. For $t \ge 3$, if the claim remains un-flipped, an adaptive policy selects among the remaining 20 attacks conditioned on state $\mathbf{s}_t$ and tried-action mask $\mathbf{m}_t$. We evaluate two distinct instantiations: (1)~\textbf{Hybrid REINFORCE Auditor}, optimizing expected cumulative reward to maximize discovery yield; and (2)~\textbf{Coverage-Aware Policy}, an offline greedy policy optimizing cumulative coverage on training splits under the same query budget. Crucially, on each fold, train-split blind spots are defined strictly from training outcomes: claims in $D_{\text{train}}$ that resist the train top-5 attacks but flip under an alternative arm. The offline greedy policy optimizes coverage strictly over $D_{\text{train}}$, without observing test-split outcomes or corpus-wide definitions.

\begin{table}[!t]
\caption{Budgeted attack selection benchmark ($K \le 5$, five seeds, mean $\pm$ sample SD, ddof=1) on held-out test splits ($N=167$ claims; 60/20/20 train/val/test). Offline policies use train outcomes only. Discovery: claims with $\ge 1$ verified flip; Blind-Spot: pooled recall across 35 held-out evaluation instances over the 34 blind-spot claims; Avg Steps: mean of per-split median steps to flip.}
\label{tab:selection}
\centering
\begingroup
\renewcommand{\arraystretch}{0.80}
\resizebox{\columnwidth}{!}{%
\begin{tabular}{lrrr}
\toprule
Method / Policy & Discovery (\%) & Blind-Spot (\%) & Avg Steps \\
\midrule
\multicolumn{4}{l}{\emph{Online Dynamic (No prior outcome knowledge)}} \\
Random-5 Selection & 40.84 $\pm$ 3.29 & 28.57 & 2.60 \\
UCB Bandit (Claim-Agnostic) & 72.69 $\pm$ 4.43 & 14.29 & 1.00 \\
\midrule
\multicolumn{4}{l}{\emph{Warm-Start Exploration Baselines}} \\
Top-2 + Random (Steps 3--5) & 66.23 $\pm$ 2.83 & 22.86 & 2.00 \\
\textbf{Hybrid Auditor (Top-2 + RL)} & \textbf{83.71 $\pm$ 2.66} & 0.00 & \textbf{1.20} \\
\midrule
\multicolumn{4}{l}{\emph{Offline Benchmark Policies (Train-split outcomes)}} \\
Pure RL Selector ($\mathbb{R}^{859}$) & 82.16 $\pm$ 2.26 & 0.00 & 1.20 \\
Static Top-5 Benchmark (ASR-Ranked) & 86.71 $\pm$ 1.30 & 0.00 & 1.40 \\
Coverage-Aware Policy (Offline Greedy) & \textbf{86.71 $\pm$ 1.30} & \textbf{51.43} & 1.40 \\
\midrule
\multicolumn{4}{l}{\emph{Empirical Reference}} \\
Oracle-22 (22-Attack Empirical Ceiling) & 88.70 & 100.00 & 22.00 \\
\bottomrule
\end{tabular}%
}
\endgroup
\end{table}
Table~\ref{tab:selection} reports multi-seed benchmark results. Key findings:
1. \textbf{Handover Yield}: The Hybrid Auditor achieves a \textbf{+17.49 pp improvement} over Top-2 + Random (66.23\%$\to$83.71\%; per-seed test scores: 80.84\%, 83.23\%, 85.03\%, 86.83\%, 82.63\% vs.\ 66.68\%, 67.15\%, 71.75\%, 65.69\%, 69.54\%, paired improvements $+14.16$, $+16.08$, $+13.28$, $+21.14$, $+13.09$ pp, paired $t$-test $p = 0.0004$; test-only claim bootstrap 95\% CI $[{+}12.80, {+}18.33]$ pp on held-out test splits), recovering \textbf{37.0 test claims/split} that resisted Top-2 anchors (185 total; mean queries 2.56 vs.\ 2.48 for Static Top-5).
2. \textbf{Average Yield vs.\ Vulnerability Coverage}: Static Top-5 achieves 86.71\% $\pm$ 1.30\% on test splits (84.62\% across corpus) by greedily exploiting evidence poisoning, but achieves \textbf{0/34 blind-spot recall by construction}. Return-maximizing RL similarly yields \textbf{0\% recall across 35 seed-level blind-spot evaluation instances} (vs.\ 22.86\% for Top-2 + Random; Fisher's exact $p = 0.005$), because \texttt{Fact2Fiction}, \texttt{FactMixing}, and \texttt{ClaimRewrite} drive over 92\% of RL actions.
3. \textbf{Coverage-Aware Auditing}: Optimizing for~blind-spot coverage on training splits replaces \texttt{Fact2Fiction} with \texttt{Word Jumbling} in 4/5 folds, recovering \textbf{18/35 (51.43\%) blind-spot instances} on held-out test splits while identically matching Static Top-5's 86.71\% average discovery. Because blind spots represent a rare subset (34 claims, 5--9 per fold), fold recall exhibits dispersion (44.4\%, 85.7\%, 0.0\%, 57.1\%, 80.0\%; unweighted macro $53.5\% \pm 33.7\%$). Fold 3 (0\%) contained only non-jumbling claims, demonstrating that multi-fold pooling is essential and that optimizing for expected yield blinds auditors to long-tail failures. The 51.43\% blind-spot recovery should be interpreted with the Word Jumbling limitation in mind: this attack preserves lexical content ($J=1.0$) but can produce awkward constituent order, so the result does not represent fully fluency-validated coverage. Among the 20 non-jumbling blind-spot claims, 15 are vulnerable to surface perturbations; 8 are strictly surface-only after excluding both Word Jumbling and syntactic omission. Recovering these cases motivates further exploration of morphological and syntactic omission attacks.
4. \textbf{Ablation \& MDP Implementation}: Masking semantic features ($\mathbf{z}_c \in \mathbb{R}^{768}, \mathbf{v} \in \{0, 1\}^3$, leaving tried-action mask $\mathbf{m}_t \in \{0, 1\}^{22}$ and outcome history $\mathbf{h}_t \in \{0, 1\}^{66}$, $\mathbf{s}_t \in \mathbb{R}^{859}$) shifts Pure RL test discovery from 82.16\% $\pm$ 2.26\% to 79.76\% $\pm$ 2.45\% ($-$2.40 pp; and Hybrid Auditor discovery from 83.71\% $\pm$ 2.66\% to 81.32\% $\pm$ 2.51\%, $-$2.39 pp), confirming that empirical priors and action masking drive primary search efficiency while semantic embeddings provide complementary conditioning. The policy $\pi_\theta$ ($\text{Linear}(859, 256) \to \text{ReLU} \to \text{Linear}(256, 22)$) and baseline $V_\phi$ are trained via REINFORCE with query penalty $-0.05$, flip reward $+1.0$, POS bonus $+0.5$, $\gamma=0.99$, and entropy bonus $\beta=0.05$ across seeds 42--46.

\section{Cross-Architecture Generalization \& Transferability (RQ2)}
\label{sec:transfer}

While \texttt{gpt-4o-mini} serves as our standardized primary verifier for exhaustive 22-attack profiling across all 24,640 instances, a critical scientific question is whether the uncovered vulnerabilities represent generalizable fact-checking failure modes or quirks of a single proprietary decoder. To evaluate cross-family transferability, 966 completed verification pairs representing successful flips on \texttt{gpt-4o-mini} were audited on open-weight \texttt{Meta-Llama-3-70B-Instruct} (Table~\ref{tab:transfer}). Across all 6 audited attack arms, \textbf{707 of 966 flips reproduced (73.19\% overall transfer)}. For arms with $>200$ flips, 200 items were randomly sampled without replacement; for smaller arms (Claim Rewrite $n=164$, Retrieval Noise $n=2$), all available flips were evaluated.

\begin{table}[!t]
\caption{Cross-model transfer audit on Llama 3 70B Instruct across 966 completed pairs (evaluating flips identified by \texttt{gpt-4o-mini}). Reprod.\ Rate = Confirmed / Audited.}
\label{tab:transfer}
\centering
\begingroup
\renewcommand{\arraystretch}{0.84}
\resizebox{\columnwidth}{!}{%
\setlength{\tabcolsep}{3.5pt}
\begin{tabular}{lrrr}
\toprule
Attack Mechanism & Audited & Confirmed & Reprod.\ (\%) \\
\midrule
Agentic fictional evidence (Fact2Fiction) & 200 & 200 & \textbf{100.00} \\
Poisoned evidence addition (AdvAdd) & 200 & 199 & \textbf{99.50} \\
Contextualized evidence replace & 200 & 179 & \textbf{89.50} \\
Masked-token claim rewrite & 164 & 121 & 73.78 \\
Imperceptible retrieval noise & 2 & 1 & 50.00 \\
Fact mixing & 200 & 7 & \textbf{3.50} \\
\midrule
\textbf{Total (6 Attack Arms)} & \textbf{966} & \textbf{707} & \textbf{73.19} \\
\bottomrule
\end{tabular}%
}
\endgroup
\end{table}

\textbf{Cross-Architecture Reproduction on Audited Evidence-Manipulation Attacks.} Across the three audited evidence arms spanning Threat Model B (\texttt{AdvAdd}, \texttt{Fact2Fiction}) and Threat Model C (\texttt{ContextReplace}), 578 of 600 flips reproduced on Llama 3 70B (\textbf{96.33\% reproduction rate}), demonstrating that evidence-manipulation attacks exhibit strong cross-family transfer in our audited evaluation. Across all 6 audited arms, 707 of 966 flips reproduced (73.19\%).

\textbf{Fact Mixing Observed as Substantially Verifier-Dependent (3.50\%).} In sharp contrast, Fact Mixing exhibits acute verifier-specific sensitivity: while achieving 55.81\% conditional ASR on \texttt{gpt-4o-mini}, only 3.50\% (7/200) of flips reproduce on Llama 3 70B. In reciprocal evaluations where Llama 3 70B generated attacks ($n=25$ claims, 100 evaluations; 58.0\% overall ASR), Fact Mixing flipped \texttt{gpt-4o-mini} on 80.0\% of claims. Thus, Fact Mixing's ASR varies from 55.8\% (\texttt{gpt-4o-mini} primary audit) to 3.5\% (Llama 3 transfer) and 80.0\% (reciprocal audit against \texttt{gpt-4o-mini}), revealing that sensitivity to multi-claim semantic blending does not transfer directly across architectures, and confirming that fixed empirical rankings are inherently verifier-dependent.

\textbf{Evaluation Across Decoder and Encoder Architectures.} Beyond generative decoders, we evaluated zero-shot multilingual NLI (\texttt{mDeBERTa-v3-base-xnli}, \citealt{laurer2023building}) and a task-specialized Indic encoder (\texttt{IndicBERTv2-MLM-only}, \citealt{doddapaneni2023indic}) fine-tuned on 2,801 authentic Hindi fact-checking pairs. Training pairs were sourced from independent journalistic fact-checking registries with strict MinHash claim and event deduplication (Jaccard similarity $< 0.30$) ensuring zero overlap with the 167 held-out test claims. On clean claims, \texttt{mDeBERTa-v3} achieves 52.50\% zero-shot accuracy ($N=1,120$), whereas fine-tuned \texttt{IndicBERTv2} achieves \textbf{95.81\% clean accuracy} on held-out test claims (160/167; 98.92\% SUP, 90.32\% REF, 100.0\% NEI). Under attacks conditioned on clean-correct instances, both encoders exhibit acute evidence vulnerability: for \texttt{mDeBERTa-v3} ($N=588$ clean-correct claims), ContextReplace achieves \textbf{77.04\% ASR} (453/588) and \texttt{AdvAdd} achieves \textbf{71.60\% ASR} (421/588); for fine-tuned \texttt{IndicBERTv2} ($N=160$ clean-correct test claims), \texttt{Fact2Fiction} achieves \textbf{48.75\% ASR} (78/160), \texttt{AdvAdd} achieves \textbf{43.75\% ASR} (70/160; 64.3\% \texttt{SUP}$\to$\texttt{REF}), and ContextReplace achieves \textbf{29.38\% ASR} (47/160).

Cross-paradigm evaluation reveals \textbf{contrasting failure modes (Overwrite vs.\ Epistemic Collapse)}: for the primary verifier, \texttt{gpt-4o-mini}, 88.2\% of \texttt{AdvAdd} flips transition from gold SUP to attacked REF, with fewer than 1\% transitioning to NEI. In contrast, zero-shot \texttt{mDeBERTa-v3} exhibits epistemic collapse under \texttt{AdvAdd}: 68.9\% of its flips transition to NEI (182 SUP$\to$NEI, 108 REF$\to$NEI; 13 SUP$\to$REF), defaulting to Neutral on premise tension. Fine-tuned \texttt{IndicBERTv2} forces direct overwrite (SUP$\to$REF in 45/70 \texttt{AdvAdd} flips), as classification heads suppress NEI. Furthermore, while decoder LLMs exhibit subword resilience to character noise ($<$3.5\% ASR), dense encoders are acutely fragile to surface perturbations: \texttt{mDeBERTa-v3} suffers \textbf{78.06\% ASR} under Character Insertion (97.2\% flips to NEI), while \texttt{IndicBERTv2} suffers \textbf{28.75\% ASR} on Character Insertion and \textbf{31.25\% ASR} on Word Jumbling, consistent with character noise fragmenting Devanagari subwords.

\section{Conclusion}
\label{sec:conclusion}
We presented HAFT, an adversarial benchmark for Hindi automated fact-checking spanning 18,304 clean-verified evaluations (24,640 total across 1,120 claims and 22 attacks). Evidence poisoning achieves 41--48\% unconditional ASR on the primary decoder verifier. Among sampled evidence-attack failures identified by \texttt{gpt-4o-mini}, 578 of 600 reproduced on Llama 3 70B (96.33\%). Across all six audited attack arms, 707 of 966 sampled failures reproduced (73.19\%). On the primary decoder verifier, character noise produces $<$3.5\% ASR, whereas dense encoders exhibit substantially greater sensitivity to some surface perturbations. Decoder LLMs exhibit uncalibrated evidence overwrite, while dense encoders exhibit epistemic collapse (\texttt{mDeBERTa-v3}) or direct overwrite (\texttt{IndicBERTv2}) under poisoning, alongside acute fragility to surface perturbations. Under constrained budgets ($K \le 5$), the Hybrid Auditor achieves 83.71\% $\pm$ 2.66\% discovery (+17.49 pp over Top-2 + Random; 37.0 rescues/split). A coverage-aware policy recovered 51.43\% of blind spots while retaining Static Top-5-level discovery (86.71\% $\pm$ 1.30\%). Overall, HAFT demonstrates that attack selection maximizing average discovery does not necessarily maximize vulnerability coverage: while fixed rankings provide strong average yield, defensive auditing requires explicit coverage-aware exploration to discover the long tail of claim-specific vulnerabilities.

\section*{Limitations}
\label{sec:limitations}

We highlight six principal limitations of our work:
\begin{enumerate}[leftmargin=1.4em,itemsep=1pt,topsep=2pt]
\item \textbf{Catalogue Scope}: While HAFT benchmarks 22 representative attack mechanisms across 18,304 clean-verified evaluations (24,640 total evaluations across 1,120 claims), 31 candidate mechanisms from our broader 53-attack taxonomy remain unmeasured prioritization hypotheses for future expansions.
\item \textbf{Single-Passage Evidence}: Current benchmark instances evaluate single evidence passages; multi-hop temporal reasoning and multi-document synthesis are left for future extensions.
\item \textbf{Specialized Task Fine-Tuning}: While we evaluate across decoder LLMs (\texttt{gpt-4o-mini}, \texttt{Llama-3-70B}), cross-lingual encoders (\texttt{mDeBERTa-v3}), and an in-domain fine-tuned Indic encoder (\texttt{IndicBERTv2}), scaling fine-tuning across broader open-weights families (e.g., Qwen 2.5, Llama 3.3, MuRIL) and multi-hop retrieval is an important avenue for future benchmarking.
\item \textbf{Evidence-Context Injection Simulation}: Threat Model B simulates retrieval poisoning by concatenating fabricated passages into the evidence context rather than poisoning live web indices or dynamic search rerankers.
\item \textbf{Grammaticality of Word Jumbling}: While Hindi's postpositional case system allows constituent flexibility, arbitrary token permutation can impair natural readability ($J=1.0$).
\item \textbf{Offline Reference vs.\ Adaptive Deployment}: Static empirical policies provide a useful offline reference for average discovery under an established attack distribution, but serve as offline references rather than claim-adaptive deployment policies because they do not condition attack selection on claim-level audit history or transfer directly across distinct verifier architectures.
\end{enumerate}

\section*{Ethics and Broader Impact}
\label{sec:ethics}

Automated fact-checking systems are deployed to protect public information integrity in linguistically diverse regions. By formalizing a 77.27\% reduction in verifier query budgets, HAFT enables resource-constrained civic organizations to stress-test their models. Sourced claims originate from public journalistic archives without private personal identifiable information (PII). All evaluations were executed offline against frozen benchmark claims; no live web endpoints or third-party fact-checking portals were queried or poisoned. Fabricated Devanagari passages are dual-use; to prevent weaponization, benchmark datasets and code will be released under CC BY-NC-SA 4.0 strictly for non-commercial defensive research upon publication, with an anonymized supplementary package (including evaluation scripts, quality gating rules, confusion tables, and accepted/rejected examples) provided for review, and raw evaluation traces available upon request. Native annotators involved in quality calibration received fair above-market compensation and gave informed consent.

\bibliography{custom}

\end{document}
